\scanpage{139}
\mspage{- 132 -}
% archive sheet number: 138
\hfill n\textsuperscript{o} 316

terme constant, et utiliser le fait que la sous-algèbre des $X\handwritten{\in}\mathit{gl}(n,K)$ telles que $ {}^{t}XR+RX=0$ est de dimension $n(n-1)/2$ ).

c) Définir de même le \uline{groupe symplectique formel} à $2n$ variables sur un corps (commutatif) $K$ quelconque, et montrer que son algèbre de Lie s'identifie à la sous-algèbre de $\mathit{gl}(2n,K)$ formée des matrices $X$ telles que $ {}^{t}XA+AX=0$, où $A=\handwritten{\Bigl(}\begin{matrix}0&I_n\\-I_n&0\end{matrix}\handwritten{\Bigr)}$

27) - Soient $K$ un corps de caractéristique $p>0$, $G$ le groupe formel de dimension 2 défini par $f_1(x,y)=X_1+Y_1+X_1Y_1$, $f_2(x,y)=X_2+Y_2(1+X_1^p)$. Montrer que $G$ n'est pas commutatif, mais que son algèbre de Lie est commutative.
